\worksheettitle
  {Условия между цифрами}
  {\modulemark{M07} Версия учителя}

\begin{teacherbox}[Паспорт модуля]
  \textbf{Результат:} ученик переводит условие между цифрами в математическую
  запись, ведёт упорядоченный перебор, находит все подходящие двузначные числа
  и обосновывает полноту списка.

  \textbf{Предварительно:} ученик различает цифру и её значение, понимает
  разряды и умеет вести перебор по первой цифре (M01, M05).

  Модуль не ограничен временем. Объём практики выбирается по устойчивости
  способа, а переход определяется выполнением контрольной точки.
\end{teacherbox}

\begin{teacherbox}[Структура ученического маршрута]
  В листе явно разделены результат, стартовая задача, разобранный пример,
  алгоритм, практика с опорой, обязательная самостоятельная практика,
  дополнительные задания, разбор ошибок, контроль, самооценка и дальнейший
  маршрут.
\end{teacherbox}

\begin{teacherbox}[Смысловой центр]
  Модуль начинается не с обозначений, а с конкретной задачи о коде
  электронного замка. Главный вопрос для ученика: «Как найти все возможные
  коды и доказать, что ни один не пропущен?» Равенство и таблица вводятся как
  инструменты ответа на этот вопрос, а не как самостоятельная цель.
\end{teacherbox}

\begin{teacherbox}[Как провести стартовую задачу]
  Сначала примите два-три найденных учеником примера, но спросите: «Это все
  коды или только некоторые? Как это проверить?» Не вводите обозначения
  \(a\) и \(b\) до появления потребности упорядочить поиск. Полный ответ:
  \(13,24,35,46,57,68,79\).
\end{teacherbox}

\begin{checkpointbox}[Если нужна проверка опоры M05]
  Предложите отдельно составить из цифр \(1,3,5\) все двузначные числа без
  повторений. Ответ:
  \(13,15,31,35,51,53\). Полнота: по очереди выбрана каждая из трёх цифр
  десятков; для неё перебраны две оставшиеся цифры единиц. Если список случаен
  или неполон, сначала используйте опору M05-R.
\end{checkpointbox}

\section{От сюжета к математической модели}

Связка с M05: раньше ученик выбирал обе цифры, теперь условие определяет
вторую цифру после выбора первой.

Для двузначного числа \(10a+b\): \(1\leq a\leq9\), \(0\leq b\leq9\).
В первой модели \(b=a+2\): при \(a=1\) получаем \(b=3\) и число \(13\), при
\(a=2\) --- \(b=4\) и число \(24\).

После разбора стартовой задачи можно проверить направление зависимости:
\begin{enumerate}
  \item \(b=a+4\);
  \item \(a=b+3\), равносильно \(b=a-3\);
  \item \(b=a-5\), равносильно \(a=b+5\).
\end{enumerate}

\begin{teacherbox}[Два решения перед началом перебора]
  Сначала ученик переводит фразу в равенство. Затем выбирает цифру, которую
  удобно изменять по порядку: для \(b=f(a)\) обычно перебирается \(a\), для
  \(a=f(b)\) --- \(b\). Не закрепляйте правило «всегда начинаем с \(a\)».
\end{teacherbox}

\section{Полный список: ответы}

Для \(b=a+2\):
\[
  13,24,35,46,57,68,79.
\]
Значения \(a+2=10,11\) являются кандидатами, но не становятся цифрой \(b\).
Последняя допустимая пара --- \((7,9)\), первая недопустимая --- \((8,10)\).

Для условия «цифра единиц на \(3\) меньше цифры десятков»:
\[
  b=a-3, \qquad 30,41,52,63,74,85,96.
\]
Перебор начинаем с \(a=3\): при \(a=1\) и \(a=2\) из цифры десятков нельзя
вычесть \(3\) и получить цифру единиц в изученных числах. Первая пара ---
\((3,0)\).

\begin{teacherbox}[Четыре обязательных вопроса]
  Как записано условие? Какую цифру удобнее перебирать? Как вычисляется
  вторая цифра? Какие значения перебираемой цифры проверены, чтобы доказать
  полноту?
\end{teacherbox}

\section{Кратность и сумма: ответы}

Для \(b=2a\): \(12,24,36,48\). Последняя допустимая пара --- \((4,8)\),
первая недопустимая --- \((5,10)\).

Для \(a=2b\) удобно перебирать \(b=1,2,3,4\): \(21,42,63,84\). При \(b=0\)
получается \(a=0\), а цифра десятков двузначного числа не может быть нулём.

Для \(a+b=7\): \(16,25,34,43,52,61,70\). Число \(70\) входит в список,
потому что \(0\) является допустимой цифрой единиц.

Математические записи:
\begin{enumerate}
  \item \(b=3a\);
  \item \(a=2b\);
  \item \(a+b=9\);
  \item \(a+b=12\).
\end{enumerate}

\section{Самостоятельная практика: ответы}

\begin{enumerate}
  \item \(b=a+1\): \(12,23,34,45,56,67,78,89\). Последняя допустимая
    пара \((8,9)\), первая недопустимая \((9,10)\).
  \item \(b=a+4\): \(15,26,37,48,59\). Последняя допустимая пара \((5,9)\),
    первая недопустимая \((6,10)\).
  \item \(b=a-2\): \(20,31,42,53,64,75,86,97\). Перебор начинаем с
    \(a=2\), потому что при \(a=1\) из цифры десятков нельзя вычесть (2) и
    получить цифру единиц в изученных числах.
  \item \(a=b+4\), или \(b=a-4\): \(40,51,62,73,84,95\). Всего 6 чисел.
  \item \(b=3a\): \(13,26,39\). Последняя допустимая пара \((3,9)\), первая
    недопустимая \((4,12)\).
  \item \(a=2b\): перебираем \(b\); получаем \(21,42,63,84\). При \(b=0\)
    получается \(a=0\), а цифра десятков не может быть нулём.
  \item \(a+b=5\): \(14,23,32,41,50\). Всего 5 чисел.
  \item \(a+b=10\): \(19,28,37,46,55,64,73,82,91\). Крайние кандидаты
    \(a=0\) и \(b=10\) недопустимы: \(a\) не может быть нулём, а \(b\) не
    может быть двузначным значением.
\end{enumerate}

\begin{teacherbox}[Дозирование по результату]
  \begin{itemize}
    \item Если направление зависимости неустойчиво, остановитесь после
      заданий 1--4 и используйте M07-R.
    \item Если списки верны, но неполны, требуйте назвать все проверенные
      значения выбранной цифры; граничная пара сама по себе не заменяет
      доказательство полноты.
    \item Если способ устойчив во всех трёх типах условий, переходите к
      смешанным заданиям и M07\(\star\).
  \end{itemize}
\end{teacherbox}

\begin{teacherbox}[Как использовать сюжетные формулировки]
  Сюжет нужен для постановки вопроса и переноса способа, но не для каждого
  однотипного упражнения. В листе контекст возвращается в задании о шкафчиках
  и в части задач с дополнительным условием; остальные задания оставлены
  короткими для чистой отработки алгоритма.
\end{teacherbox}

\section{Ошибки: ответы и причины}

\begin{enumerate}
  \item Направление обращено. Верно \(b=a-2\); список:
    \(20,31,42,53,64,75,86,97\).
  \item \(10\) не является цифрой. Для \(b=a+4\) список:
    \(15,26,37,48,59\).
  \item Ноль допустим в разряде единиц. Полный список:
    \(15,24,33,42,51,60\).
  \item Перепутаны роли \(a\) и \(b\). Для \(a=2b\): \(21,42,63,84\).
  \item Перебор остановлен раньше времени. Потеряны \(47,58,69\); после
    пары \((6,9)\) получается первая недопустимая пара \((7,10)\).
  \item Запись \(4\) не является двузначным числом. Полный список:
    \(13,22,31,40\).
\end{enumerate}

\begin{teacherbox}[Диагностируйте причину, а не только ответ]
  \begin{itemize}
    \item неверный знак или переставленные переменные --- обращение зависимости;
    \item вычисления начинаются раньше первого допустимого значения или дают
      число больше \(9\) --- не удерживаются границы цифры;
    \item пропущено число с единицами \(0\) --- ноль ошибочно исключён;
    \item несколько верных примеров без полного списка --- нет упорядоченного перебора;
    \item включено однозначное число --- не удерживается ограничение \(a\geq1\).
  \end{itemize}
\end{teacherbox}

\section{Задачи с дополнительным условием: ответы}

\begin{enumerate}
  \item Сумма \(8\) и \(a>b\): \(53,62,71,80\).
  \item \(b=a+2\) и число больше \(40\): \(46,57,68,79\).
  \item \(a=2b\) и число меньше \(80\): \(21,42,63\).
  \item Сумма \(11\) и \(50<n<90\): \(56,65,74,83\).
  \item \(b=a-3\) и число нечётное: \(41,63,85\).
\end{enumerate}

Собственное равенство оценивается по четырём признакам: точная фраза,
соответствующая математическая запись, полный список и обоснование полноты.

\section{Контрольная точка}

\begin{enumerate}
  \item \(a=b+4\), или \(b=a-4\): удобно перебирать \(b=0,\ldots,5\);
    список \(40,51,62,73,84,95\).
  \item \(b=2a\): перебираем \(a=1,\ldots,4\); \(12,24,36,48\).
  \item \(a+b=8\): перебираем \(a=1,\ldots,8\);
    \(17,26,35,44,53,62,71,80\).
  \item Больше \(50\): \(53,62,71,80\).
  \item Полнота подтверждается проверкой всех допустимых значений выбранной
    цифры и сохранением каждой пары, в которой обе цифры допустимы.
\end{enumerate}

\begin{resultbox}[Решение о маршруте]
  \begin{itemize}
    \item \textbf{Результат M07 достигнут:} во всех трёх типах условий верны
      математическая запись, выбор перебираемой цифры и полный список;
      полнота объяснена через проверенные значения.
    \item \textbf{M07-R:} обращена зависимость, включены недопустимые значения
      цифры, пропущен ноль или список построен случайно.
    \item \textbf{Дополнительная практика:} способ понят, но в одном из трёх
      типов условий список остаётся неполным. После закрепления можно
      назначить M07-H.
    \item \textbf{M07\(\star\):} базовый результат устойчив; ученик готов
      считать допустимые значения без выписывания всех чисел.
  \end{itemize}
\end{resultbox}

\begin{teacherbox}[После коррекции]
  Дайте новое условие: «цифра единиц на \(4\) меньше цифры десятков».
  Ожидаются \(b=a-4\), осознанный выбор \(a\) или \(b\), список
  \(40,51,62,73,84,95\) и объяснение полноты через все проверенные значения
  выбранной цифры.
\end{teacherbox}

\newpage
\section{Ответы к M07-R и M07\(\star\)}

\begin{teacherbox}[Ответы M07-R]
  Для условия \(b=a+4\) кандидаты проверяются при
  \(a=1,2,\ldots,9\). Допустимы первые пять пар; полный список:
  \(15,26,37,48,59\). Значения \(10,11,12,13\) не являются цифрами \(b\).

  Для \(a=2b\) допустимы \(b=1,2,3,4\); полный список:
  \(21,42,63,84\). При \(b=0\) цифра десятков равна нулю, а при \(b\geq5\)
  кандидат на значение \(a\) не меньше \(10\).

  Повторная проверка: \(b=a-2\), удобно перебирать
  \(a=2,3,\ldots,9\); список \(20,31,42,53,64,75,86,97\).
\end{teacherbox}

\begin{teacherbox}[Ответы M07\(\star\): разность цифр]
  \begin{itemize}
    \item \(b=a+1\): \(a\) от 1 до 8, всего 8 чисел;
    \item \(b=a+2\): \(a\) от 1 до 7, всего 7 чисел;
    \item \(b=a+4\): \(a\) от 1 до 5, всего 5 чисел;
    \item \(b=a-1\): \(a\) от 1 до 9, всего 9 чисел;
    \item \(b=a-3\): \(a\) от 3 до 9, всего 7 чисел;
    \item \(b=a-6\): \(a\) от 6 до 9, всего 4 числа.
  \end{itemize}
  Допустимая формулировка закономерности: чем больше прибавляем к цифре
  десятков, тем раньше кандидат на цифру единиц становится больше 9 и тем
  меньше остаётся чисел. При вычитании ноль допустим как цифра единиц, поэтому
  для \(b=a-1\) подходят все девять значений \(a\).
\end{teacherbox}

\begin{teacherbox}[Ответы M07\(\star\): сумма цифр]
  \begin{itemize}
    \item суммы \(4,7,9\) дают соответственно \(4,7,9\) чисел;
    \item суммы \(10,11,12,15,18\) дают соответственно \(9,8,7,4,1\)
      вариант;
    \item после суммы 9 при каждом увеличении суммы на 1 один крайний вариант
      становится недопустимым.
  \end{itemize}

  Контроль: для \(b=a+3\) допустимы \(a=1,\ldots,6\), всего 6 чисел; для
  \(b=a-4\) допустимы \(a=4,\ldots,9\), всего 6 чисел; при сумме цифр 13
  допустимы \(a=4,\ldots,9\), всего 6 чисел.
\end{teacherbox}

\newpage
\begin{teacherbox}[Необязательное обобщение с учителем]
  Общую формулу не требуйте от всего класса. Если ученик сам готов к
  обобщению, можно обозначить постоянную разность одной новой буквой и после
  конкретных таблиц получить правило для количества. Аналогично для суммы:
  сначала формулируется словесная закономерность, и только затем, по желанию,
  вводится обозначение суммы. Не используйте обе новые буквы одновременно.
\end{teacherbox}

\begin{teacherbox}[Критерии собственных условий и объяснения полноты]
  Собственное условие принимается, если оно однозначно сформулировано словами,
  верно записано равенством, даёт от трёх до шести двузначных чисел и список
  соответствует условию.

  Допустимые объяснения полноты могут звучать по-разному. Обязательная мысль:
  ученик по порядку проверил все возможные значения выбранной цифры, для
  каждого вычислил вторую цифру и сохранил ровно те пары, в которых цифра
  десятков лежит от 1 до 9, а цифра единиц --- от 0 до 9. Для M07\(\star\)
  достаточно назвать весь допустимый промежуток значений и его длину.
\end{teacherbox}
